The capstone is short because you already know all the commands --- this
tutorial is about turning them into automation: script files, one-liners, exit
codes a shell can branch on, the save/load round trip, and the trajectory CSV.

\subsection*{Build once, save, prove the round trip}

\begin{console}
\muted{# QtRocket CLI - type 'help' for commands, 'quit' to exit.}
\uin{loaddb estes_small.qmd}
\ok loaddb: 13 new motors from estes_small.qmd (13 total)
\uin{newdesign NoseCone name=nose baseRadius=0.0098 length=0.10 density=700}
\ok newdesign: root NoseCone id=2
\uin{addpart nose BodyTube name=body innerRadius=0.0091 outerRadius=0.0098 length=0.25 density=700}
\ok addpart: BodyTube id=3 under 2
\uin{addpart body FinSet name=fins finCount=3 rootChord=0.04 tipChord=0.02 span=0.03 sweep=0.02 thickness=0.003 bodyRadius=0.0098 density=700 seat=OnSurface parentStation=0 childStation=0}
\ok addpart: FinSet id=4 under 3
\uin{setmotor C6}
\ok setmotor: C6
\uin{setdrag 0.75}
\ok setdrag: 0.75
\uin{setarea 0.000302}
\ok setarea: 0.000302 m^2
\uin{checkdesign}
\ok checkdesign: 4 parts, contiguous, no overlaps
\uin{savedesign workhorse.qrd}
\ok savedesign: workhorse.qrd
\end{console}


\begin{console}
\uin{cleardesign}
\ok cleardesign: restored the placeholder body
\uin{loaddesign workhorse.qrd}
\ok loaddesign: workhorse.qrd (motor C6)
\uin{checkdesign}
\ok checkdesign: 4 parts, contiguous, no overlaps
\uin{status}
\ok status:
  motor      = C6
  dry_mass   = 1 kg
  drag_coeff = 0.75
  ref_area   = 0.000302 m^2
  velocity   = 0 m/s
  angle      = 0 deg (from vertical)
  atmosphere = Constant Atmosphere
  gravity    = Constant Gravity
  integrator = Runge-Kutta 4th Order
  database   = 13 motors
\uin{launch}
\ok launch: 2041 steps, t_final=20.400s
  apogee    = 436.562 m @ t=7.990s
  max_speed = 143.732 m/s @ t=1.820s
  downrange = 0.000 m
  landing   = t=20.400s, pos=(0.000, 0.000, -0.436)
CSV /tmp/flight1/qtrocket_run.csv
\end{console}


\repl{cleardesign} wipes the session; \repl{loaddesign} restores everything
from the \code{.qrd} --- tree, placement, drag configuration, and the motor
re-resolved \emph{by name} from the loaded database. The reloaded rocket
passes the same \repl{checkdesign} and flies the same flight. That is the
round-trip contract, and it is what makes design files trustworthy carriers of
work between sessions (and between the CLI, the GUI, and the visualizer).

\subsection*{Sampling the trajectory}

\begin{console}
\uin{states 200}
\ok states: 2041 total, stride=200
t,x,y,z,vx,vy,vz,mass,cg_x,cg_y,cg_z,Ixx,Iyy,Izz
0,0,0,0,0,0,0,0.0440836721,0,0,-0.131975423,0.000430023519,0.000430023519,5.27738786e-06
2,0,0,185.449796,0,0,129.355482,0.0332836721,0,0,-0.150463003,0.000378959825,0.000378959825,4.83998786e-06
4,0,0,347.944129,0,0,50.7177043,0.0332836721,0,0,-0.150463003,0.000378959825,0.000378959825,4.83998786e-06
6,0,0,416.493841,0,0,20.6988221,0.0332836721,0,0,-0.150463003,0.000378959825,0.000378959825,4.83998786e-06
8,0,0,436.561845,0,0,-0.0492822118,0.0332836721,0,0,-0.150463003,0.000378959825,0.000378959825,4.83998786e-06
10,0,0,417.367307,0,0,-18.6519424,0.0332836721,0,0,-0.150463003,0.000378959825,0.000378959825,4.83998786e-06
12,0,0,365.250398,0,0,-32.4712673,0.0332836721,0,0,-0.150463003,0.000378959825,0.000378959825,4.83998786e-06
14,0,0,291.30428,0,0,-40.6424912,0.0332836721,0,0,-0.150463003,0.000378959825,0.000378959825,4.83998786e-06
16,0,0,205.327326,0,0,-44.8380887,0.0332836721,0,0,-0.150463003,0.000378959825,0.000378959825,4.83998786e-06
18,0,0,113.399757,0,0,-46.8368018,0.0332836721,0,0,-0.150463003,0.000378959825,0.000378959825,4.83998786e-06
20,0,0,18.6884925,0,0,-47.7549149,0.0332836721,0,0,-0.150463003,0.000378959825,0.000378959825,4.83998786e-06
\uin{save workhorse_flight.csv}
\ok save: /tmp/flight1/workhorse_flight.csv
\uin{quit}
\ok bye
\end{console}


\repl{states} prints every $n$-th state sample to the terminal --- the same
full CSV rows \repl{save} writes, all fourteen columns --- for a quick look
without leaving the session; \repl{save} writes the complete series to a
file. The columns are self-describing:

\begin{minted}{text}
t,x,y,z,vx,vy,vz,mass,cg_x,cg_y,cg_z,Ixx,Iyy,Izz
\end{minted}

--- time, position, velocity, plus the composite mass properties at every
step (mass falls during the burn; the inertia columns are what a future 6-DOF
mode will consume). Every plot in this guide is drawn from these files.

\subsection*{One-liners and exit codes}

Real captures.

\begin{minted}{text}
$ qtrocket-cli -c "loaddb estes_small.qmd" ; echo $?
# QtRocket CLI - type 'help' for commands, 'quit' to exit.
OK loaddb: 13 new motors from estes_small.qmd (13 total)
0
$ qtrocket-cli -c "launch" ; echo $?
# QtRocket CLI - type 'help' for commands, 'quit' to exit.
ERR launch: no motor set (use loadmotors + setmotor first)
1
$ qtrocket-cli tutorial5-missing.cli ; echo $?
error: cannot open script file 'tutorial5-missing.cli'
2
\end{minted}

The middle example is the process-isolation warning from
\cref{sec:invocation}: each \code{-c} is a fresh session, so the
launch never saw a motor. The exit-code --- 0 all-OK, 1 any
\code{ERR}, 2 bad invocation --- makes the CLI usable in
scripts and CI pipelines:

\begin{minted}{text}
qtrocket-cli nightly_designs.cli || echo "a design regressed"
\end{minted}

\begin{tip}[What you learned --- and where to go next]
Scripts are just sessions in a file; \code{.qrd} files carry designs across
sessions and tools; the CSV carries flights into your plotting tool of
choice; and the exit code turns any script into a regression guard ---
Tutorial~3's deliberate failures exited 1 for exactly this reason. From here:
the command reference (\cref{sec:reference}) is the complete vocabulary, and
the design whitepapers under \srcf{docs/} go as deep as you care to dig.
\end{tip}
